% %%%%%%%%%%%%%%%%%%%%%%%%%%%%%%%%%%%%%
% %           Header                  %
% %%%%%%%%%%%%%%%%%%%%%%%%%%%%%%%%%%%%%

% %-----------------------------------
% %           Introduction           %
% %-----------------------------------
\documentclass[letterpaper,11pt,oneside]{article}

%--------------------------------
% Packages
%--------------------------------
\usepackage[T1]{fontenc}
\usepackage[english]{babel}
\usepackage{amsmath, amssymb}
\usepackage[table]{xcolor}
\usepackage{times}
\usepackage{helvet}
\usepackage{sectsty}
\usepackage{enumitem} %bullet point list formatting
\usepackage{verbatim}  % comments

% Common left indent for multi-line derivations (used with flalign)
\newcommand{\derivindent}{\hspace{2em}}

% Format bullet point lists
\setlist[itemize]{itemsep=0pt, topsep=0pt, parsep=0pt, partopsep=0pt}

% Use sans serif font for sections
\allsectionsfont{\sffamily}

% Prevent footnotes to be broken across pages
\interfootnotelinepenalty=10000

\usepackage[labelfont={bf,sf}]{caption}
\usepackage{booktabs}
\usepackage{hyperref}
\usepackage{multirow}
\usepackage{pdflscape}
\usepackage[onehalfspacing]{setspace}
\usepackage{orcidlink}
\usepackage{bbold}
\usepackage{enumitem}
\usepackage{subcaption} %for subfigures
\usepackage{float} %add placement option for figures and tables
\usepackage{bm} %for bold math symbols
\usepackage{placeins} %for float barrier

% Tables
\usepackage{colortbl}
\usepackage{tabularx}
\usepackage{threeparttable}
\usepackage{xltabular}
\usepackage{longtable}
\usepackage{makecell} %complex table formatting

%--------------------------------
% Colors
%--------------------------------
\definecolor{pigment}{rgb}{0.2,0.2,0.6}
\definecolor{lightgray}{gray}{0.9}
\definecolor{lightgrey2}{RGB}{235,235,235}

\hypersetup{
    colorlinks=true,
    urlcolor=pigment,
    linkcolor=pigment,
    citecolor=pigment
}

% Revision color
\definecolor{revcol}{rgb}{0,0,0}
\newcommand{\rev}[1]{\textcolor{revcol}{#1}}

%--------------------------------
% References
%--------------------------------
\usepackage[
    style=apa,
    sortcites=true,
    sorting=nyt,
    backend=biber
]{biblatex}

% Allow long URLs in the bibliography to break after letters and digits (only if needed)
\setcounter{biburlnumpenalty}{9000}
\setcounter{biburlucpenalty}{9000}
\setcounter{biburllcpenalty}{9000}
\setcounter{biburlnumpenalty}{9000}
\setcounter{biburlucpenalty}{9000}
\setcounter{biburllcpenalty}{9000}

\usepackage{nameref}
\addbibresource{BES.bib}

\usepackage{doi}
\usepackage{csquotes}

\DeclareFieldFormat{addendum}{
    \textbf{\textit{#1}}
}

%--------------------------------
% Figures and code
%--------------------------------
\usepackage{graphicx}
\usepackage{listings}

\definecolor{codeback}{HTML}{F5F5F5}
\definecolor{codecomment}{HTML}{6C757D}
\definecolor{codestring}{HTML}{E78A00}
\definecolor{codeblue}{HTML}{085AB3}
\definecolor{johnsonblue}{HTML}{132B69}

%--------------------------------
% TikZ
%--------------------------------
\usepackage{tikz}
\usepackage{fontawesome5}

\usetikzlibrary{
    shapes.geometric,
    arrows.meta,
    positioning,
    fit,
    shadows,
    decorations.pathmorphing
}

\definecolor{zenodocolor}{RGB}{0,123,191}
\definecolor{githubcolor}{RGB}{36,41,46}
\definecolor{webcolor}{RGB}{255,87,51}
\definecolor{packagecolor}{RGB}{76,175,80}
\definecolor{usercolor}{RGB}{156,39,176}
\definecolor{datacolor}{RGB}{255,193,7}

%--------------------------------
% Margins
%--------------------------------
\usepackage{geometry}

\geometry{
    left=20mm,
    right=25mm,
    top=25mm,
    bottom=25mm
}

%--------------------------------
% Title and authors and URLs
%--------------------------------
\title{
    \vspace{-3em}
    \textsf{\textbf{
        Specifying meaningful hypotheses across studies: Bayesian evidence synthesis revisited
    }}
}

\usepackage{authblk}

\renewcommand*{\Authands}{,}
\renewcommand\Affilfont{\small}

\author[1]{
    Laura~Groot \orcidlink{0009-0000-6989-0468}
}

\author[1]{
    Daniel~W.~Heck \orcidlink{0000-0002-6302-9252}
}

\affil[1]{
    Psychological Methods Lab, Department of Psychology, Philipps-Universität Marburg
}

\date{\today}

% Link to the Supplementary Material (TODO)
\newcommand{\suppurl}{\href{https://github.com/L-Groot/BES}{\nolinkurl{https://github.com/L-Groot/BES}}}

%--------------------------------
% Document
%--------------------------------
\begin{document}

\maketitle

\begin{abstract} % max: 250 words

\noindent
Bayesian evidence synthesis (BES) is a method for evaluating the empirical support for a common theory across multiple studies. Evidence is aggregated by multiplying the Bayes factors of the individual studies, yielding the product Bayes factor (PBF). The numerator of the PBF corresponds to the joint predicted hypothesis that the theory holds across different contexts, implying that the theory's predictions (i.e., a set of order constraints on the parameters of interest) are correct in each study. Since the Bayes factor is a relative measure of evidence, interpreting the PBF requires a clear understanding of the alternative hypothesis. The denominator of the PBF corresponds to a joint alternative hypothesis stating that the respective study-specific alternative hypothesis holds in each study. This joint alternative may lack both a clear substantive interpretation and a theoretical justification. Focusing on three common types of alternatives for order-constrained inference (the null, the complement, and the unconstrained alternative), we investigate how different combinations of study-specific alternative hypotheses give rise to distinct joint alternative hypotheses. In particular, multiplying Bayes factors involving study-specific complements does not test the theory against its complement (i.e., whether the theoretical predictions are false in at least one study).
We provide recommendations for specifying meaningful joint alternative hypotheses that support a clear interpretation of BES results and are useful for theory testing.

\noindent
\textsc{Keywords}: Bayesian evidence synthesis, product Bayes factor, informative hypotheses, meta-analysis, order-constrained inference

\vfill

\noindent
Correspondence concerning this article should be addressed to 
\url{laura.groot@uni-marburg.de}.

\end{abstract}

\newpage

\section{Introduction}

\subsection{Evidence synthesis for conceptual replications}

A central goal of science is to test theories that assume relationships between latent constructs. Because constructs are by definition unobservable, testing a theory requires researchers to rely on specific operationalizations \parencite{stroebeAllegedCrisisIllusion2014}. The predictions of the theory should hold irrespective of the choice of a particular operationalization. More generally, ``theoretical statements made by researchers almost without fail imply some degree of robustness to irrelevant variables'' \parencite[p. 2608]{baribaultMetastudiesRobustTests2018}. For example, dual-process theories of persuasion \parencite{chaikenHeuristicSystematicInformation1980} posit, among other things, that deeper processing of a communication increases persuasion. If dual-process theory is correct, it should not matter whether processing depth was operationalized by manipulating personal relevance \parencite{pettyPersonalInvolvementDeterminant1981} or the expectation to discuss the communication later \parencite{chaikenHeuristicSystematicInformation1980}. 

When researchers intentionally vary theoretically irrelevant aspects of a study while testing the same theory, they perform a \emph{conceptual replication} \parencite{schmidtShallWeReally2009}. Besides varying operationalizations, conceptual replications may differ in participant populations, experimental procedures, statistical analyses, or aspects of the study design, such as timing and number of measurement occasions. When a theory is supported across different constellations of theoretically irrelevant variables, our confidence in it increases because the proposed relationships generalize across contexts \parencite{schmidtShallWeReally2009, stroebeAllegedCrisisIllusion2014, lakensImprovingYourStatistical2022}. Hence, it is often of interest to synthesize evidence for or against a theory across a set of conceptual replications.

Yet the very heterogeneity that gives conceptual replications higher evidential value poses challenges for statistical analysis. For example, Bayesian sequential updating \parencite{schonbrodtSequentialHypothesisTesting2017} pools the raw data across all studies, which requires all studies to share the same data format and model structure. This makes Bayesian sequential updating suitable for evidence synthesis across a set of direct replications, but not for conceptual replications. In comparison, standard meta-analysis approaches are more flexible because aggregation occurs at the level of effect sizes. Specifically, random-effects meta-analysis assumes between-study heterogeneity and allows the true effect size to vary across studies.

However, conventional meta-analysis may not be suitable for synthesizing conceptual replications because it assumes that study-specific effect sizes are comparable \parencite{borensteinBasicIntroductionFixedeffect2010}, that is, they can be expressed as a common effect-size metric (e.g., all studies estimate Cohen's~$d$). For conceptual replications, this assumption is usually violated. First, different studies test different predictions derived from the same theory. For example, the first study may examine a simple mean difference, the second a complex interaction, and the third a certain ordering of regression coefficients. The studies thus involve different parameters---sometimes multiple---that cannot be reduced to a single effect size. Moreover, even if each study yields one effect size, it is unlikely that all studies use the same measure. Although different effect sizes can sometimes be converted into a common metric, such transformations can introduce ``strong or even untenable assumptions, resulting in estimates of effect sizes with unknown distributional properties'' \parencite[p.~1]{vanassenEndJustifiesAll2023}. Second, effect-size comparability can be violated even if all studies test the same hypothesis using identical designs and statistical methods, for example, when each study controls for a different set of covariates \parencite{vanlissaTutorialAggregatingEvidence2024}.


\subsection{Bayesian evidence synthesis (BES)}

Bayesian evidence synthesis (BES; \cite{kuiperCombiningStatisticalEvidence2013}) addresses the limitations of existing synthesis methods by accommodating between-study heterogeneity in hypotheses, study designs, and statistical methods.%
    \footnote{We use ``BES'' to refer to the evidence synthesis method of multiplying Bayes factors across studies (e.g., \cite{kevenaarBayesianEvidenceSynthesis2021, klugkistBayesianEvidenceSynthesis2025}), rather than to Bayesian meta-analytic methods in general (e.g., \cite{scheibehenneBayesianEvidenceSynthesis2016, mulderBayesFactorHypothesis2026}).}
BES requires that each study tests at least one theory-derived hypothesis and that the evidence for these study-specific hypotheses is quantified by a Bayes factor (BF). The core idea of BES is straightforward and only requires multiplication of the study-specific Bayes factors to obtain the product Bayes factor (PBF). The PBF quantifies the combined evidence for the underlying theory that is provided by all studies together. The two main assumptions of BES are that (1) all studies test the same underlying theory and (2) each study provides independent evidence, which concerns the study-specific hypotheses, the data, and the prior distributions.

BES has intuitive appeal and has been applied to a range of empirical studies, for example, to synthesize evidence on the effect of parental age on offspring mental health \parencite{zondervan-zwijnenburgParentalAgeOffspring2020}, the development of self-control during adolescence \parencite{zondervan-zwijnenburgRobustLongitudinalMulticohort2020}, and the effect of network embeddedness on trust and cooperation \parencite{kevenaarBayesianEvidenceSynthesis2021, volkerFutureMadeToday2022}. However, previous applications of BES also reveal an important blind spot. Since BES rests on computing the PBF, the interpretation of its results inherits a key property of Bayesian model selection, namely that evidence is relative. The Bayes factor does not test a hypothesis in an absolute sense but rather quantifies the evidence for one hypothesis \textit{relative to another} \parencite{tendeiroReviewIssuesNull2019}. Therefore, a correct interpretation of the PBF requires a clear specification of the competing claim against which the theory is evaluated. However, previous applications of BES have often not explicitly specified the competitor hypothesis, which increases the risk of misinterpreting the results.

% This omission has two important consequences. First, it means that the evidence for the theory of interest is often not interpreted in the relative manner as required for a Bayes factor, obscuring which ``competitor'' the theory of interest has been compared to \parencite{tendeiroReviewIssuesNull2019}. Second, the choice of competitor is not merely a minor detail: different competitor claims can lead to substantially different conclusions about the evidential status of the same theory.

\subsection{Overview}

The present paper makes four main contributions. First, we clarify which joint hypotheses are being tested when multiplying study-specific Bayes factors. Second, we show how different combinations of study-specific alternative hypotheses give rise to different joint alternative hypotheses, and demonstrate that these choices can substantially affect conclusions about the combined evidence for a theory. Third, we show that multiplying Bayes factors against the study-specific complements does not yield a test of the theory against \textit{its} complement. We call the latter the complete complement and derive the computation of the Bayes factor against it from the study-specific Bayes factors alone. Finally, we provide recommendations on specifying joint alternative hypotheses that are meaningful in the context of theory testing.



%------------------------------------------------
%           Study-specific Hypotheses           %
%------------------------------------------------

\section{Study-specific hypotheses}

\subsection{Study-specific predicted hypotheses}

In BES, one first needs to select a set of studies, each of which is regarded as an independent test of the theory of interest. For each study $t$, a study-specific predicted hypothesis is derived from the theory:
    \[
    H_1^{(t)} \quad \text{for } \, t = 1,2,\dots,T.
    \]
In mathematical psychology, many theories can be translated into informative statistical hypotheses that encode substantive expectations about the observed outcomes or model parameters \parencite{heckMultinomialModelsLinear2019, heckReviewApplicationsBayes2023}. Informative statistical hypotheses are formally defined by restricting the parameters $\boldsymbol{\theta}$ of a statistical model via a set of equality constraints (e.g., $\theta_1=\theta_2$), order constraints (e.g., $\theta_1 \leq \theta_2$), and/or interval constraints (e.g., $-0.1 \leq \theta_1 \leq 0.1$). BES assumes that the theory implies such a predicted statistical hypothesis for each study. To illustrate BES, we introduce two running examples with typical order constraints in the context of choice probabilities and equivalence testing.

\subsubsection{Example A: Choice probabilities}

Suppose we want to evaluate the evidence for the existence of precognition, defined by \textcite[][p. 410]{bemFeelingFutureExperimental2011} as ``accessing information yet to be determined in the future.'' Experiment~1 by \textcite{bemFeelingFutureExperimental2011} investigated precognition for positively valenced stimuli by showing participants two images of curtains on a display, one with a picture behind it and the other with blank space. Participants were asked to select the curtain behind which they believed the picture to be. The location of the images was randomly determined \textit{after} the participant had made their decision. In this experiment, the positively valenced stimuli were erotic images, and Bem hypothesized that ``participants would be able to identify the position of the hidden erotic picture significantly more often than chance (50\%)'' \parencite[p. 409]{bemFeelingFutureExperimental2011}.
Suppose that we want to analyze the data of three replication studies of Bem's Experiment~1. Based on the theoretical assumption that people have precognition for positive stimuli, we derive the same study-specific informative hypothesis for the hit rate $\theta^{(t)}$ in all studies:
    \[
    H_{1}^{(t)}: \theta^{(t)} > 0.50 \quad \text{for } t = 1,2,3.
    \]
The left panel of Figure~\ref{fig:study-hypotheses} illustrates the study-specific hypotheses for each of the three replication studies in pink. BES may answer the research question: \textit{How much evidence do the three replication studies jointly provide in favor of precognition for positive stimuli?}

%----Figure 1-----
\begin{figure}[!t]
\centering

\captionsetup{justification=raggedright, singlelinecheck=false}
\caption{Study-specific hypotheses for Example A (left) and Example B (right)}
\label{fig:study-hypotheses}

\vspace{12pt}

\begin{tabular}{c @{\hspace{20pt}} !{\vrule width 0.5pt} @{\hspace{20pt}} c}
\begin{subfigure}{0.45\textwidth}
    \centering
    \includegraphics[width=\linewidth]{figures/fig1_A_choice_probabilities.pdf}
\end{subfigure}
&
\begin{subfigure}{0.45\textwidth}
    \centering
    \includegraphics[width=\linewidth]{figures/fig1_B_equivalence_testing.pdf}
\end{subfigure}
\end{tabular}

\vspace{20pt}

\begin{minipage}{\textwidth}
\raggedright
\footnotesize
\textit{Note.} Predicted hypotheses (pink) and the three alternatives: null (blue), unconstrained (yellow), and complement (green).
\end{minipage}

\end{figure}

\subsubsection{Example B: Equivalence testing}

Equivalence testing aims to assess whether an observed effect is ``too small to be meaningful'' \parencite{lakensImprovingYourStatistical2022}, rather than whether it differs from the absolute null of no effect. This is realized by specifying an informative hypothesis that covers a region around an effect size of zero (i.e., no difference in means). If the observed effect falls into this region, it is considered too small to be meaningful, and statistical equivalence is inferred.

We build on Example~1 from \textcite{scheelEquivalenceTestingPsychological2017}, who reanalyzed a study on identity priming by \textcite{moonSecondaryReplicationAttempt2014}. This study sought to replicate \textcite{shihStereotypeSusceptibilityIdentity1999}, who found that Asian-American women performed better on a math test after their Asian identity had been primed, compared to a control group. \textcite{scheelEquivalenceTestingPsychological2017} tested for the absence of a meaningful difference between the identity-primed and the control group by specifying an equivalence interval for the mean difference in test scores, $\Delta_{\text{between}}$. The interval extends 6.25 percentage points above and below 0, corresponding to bounds of $\pm 0.0625$ on the proportion scale. This value was chosen to demarcate the smallest effect size of interest (SESOI)---a one-point change in grade. Therefore, \textcite{scheelEquivalenceTestingPsychological2017} specified the predicted hypothesis:
    \[
    H_{1}^{(1)}: -0.0625 < \Delta^{(1)}_{\text{between}} < 0.0625.
    \]
    
We extend the example by introducing two hypothetical conceptual replications. These also test for statistical equivalence, but each with a slightly different setup. Imagine that Study 2 has a different grading system, in which a one-point change in grade corresponds to a 10-percentage-point change in test score, resulting in a different equivalence interval:
    \[
    H_{1}^{(2)}: -0.1 < \Delta^{(2)}_{\text{between}} < 0.1.
    \]
Now suppose that Study 3 employs a within-person design in which each participant completes two math tests, one preceded by the target priming and one by the control manipulation. In this case, the parameter of interest is no longer the raw mean difference score $\Delta_\text{between}$ but the \textit{within-person} raw difference score $\Delta_\text{within}$, so the study-specific statistical hypothesis concerns a different effect size. Additionally, assume that in Study 3's grading scale, the SESOI of one grade point corresponds to 20 percentage points on the test:
    \[
    H_{1}^{(3)}: -0.2 < \Delta^{(3)}_{\text{within}} < 0.2.
    \]
The study-specific hypotheses for the three replications are illustrated in the right panel of Figure~\ref{fig:study-hypotheses} in pink. In this scenario, a question that BES could answer is: \textit{How much evidence do the three conceptual replication studies jointly provide that the identity-priming effect is too small to be meaningful?}

Example B highlights a key strength of BES, namely, that it allows researchers to synthesize studies that test different predicted hypotheses derived from the same theory. Here, the predicted hypotheses concern different effect sizes---a between-person and a within-person mean difference---and specify different equivalence intervals. In principle, BES can accommodate even more complex cases, such as studies with multiple parameters or structurally different types of constraints.


\subsection{Study-specific alternative hypotheses}

Bayes factors always compare two hypotheses against each other. In BES, in each study, the predicted hypothesis $H_1$ is tested against a competing alternative hypothesis $H_a$ by computing the study-specific Bayes factor
    \begin{equation}
    \mathrm{BF}_{1,a}^{(t)} = \frac{P(D^{(t)} \mid H_1^{(t)})}{P(D^{(t)} \mid H_a^{(t)})}, 
    \end{equation}
where $D^{(t)}$ are the data observed in study $t$.

The most common alternative for the predicted $H_1$ is the null hypothesis $H_0$, which can either be a point hypothesis that fixes the test-relevant parameters to a specific value (e.g., $\theta=0$) or constrains different parameters to be equal (e.g., $\theta_1=\theta_2$). The null generally reflects the expectation that there is no effect at all.%
    \footnote{Unlike null-hypothesis significance testing, where the null is tested against ``the alternative'' (here, $H_1$), we follow the BES literature in treating the null as one type of alternative to the predicted hypothesis.}
In turn, the unconstrained hypothesis $H_u$ imposes no constraints on the parameters at all, allowing them to vary freely over the entire parameter space. This renders the unconstrained hypothesis $H_u$ a \textit{maximally uninformative hypothesis}. As the third type of alternative, the complement hypothesis $H_c$ includes all parameter constellations that violate at least one constraint imposed by the predicted hypothesis. Substantively, $H_c$ thus captures all ways in which $H_1$ can be false.

\subsubsection{Example A: Choice probabilities}
\label{sec:example-a-choice-probabilities}

In Example A, all studies share the same set of predicted and alternative hypotheses. The left panel of Figure~\ref{fig:study-hypotheses} shows the three alternative hypotheses for each replication study. The null hypothesis, asserting the absence of precognition for positive stimuli, specifies a hit probability at chance level:
    \[
    H_{0}^{(t)}: \theta^{(t)} = 0.50 \quad \text{for } t = 1,2,3.
    \]
The unconstrained hypothesis spans the entire parameter space, in this case:
    \[
    H_{u}^{(t)}: \theta^{(t)} \in [0,1] \quad \text{for } t = 1,2,3.
    \]
Substantively, the unconstrained alternative reflects the absence of any commitment as to which choice probabilities are possible. Finally, the complement hypothesis covers all parameter values not included in $H_1^{(t)}$:
    \[
    H_{c}^{(t)}: \theta^{(t)} \leq 0.50\quad \text{for } t = 1,2,3.
    \]
Substantively, the complement hypothesis states that precognition for positive stimuli either does not exist or operates in the ``wrong'' direction, that is, in the form of avoidance rather than attraction.

\subsubsection{Example B: Equivalence testing}

In Example B, the predicted hypothesis encodes the absence of a meaningful effect by specifying an interval-null hypothesis in each study (right panel of Figure~\ref{fig:study-hypotheses}). A point null alternative is therefore typically not of interest in equivalence testing, since $H_1$ already includes the absence of an effect. Nevertheless, for conceptual clarity, Figure~\ref{fig:study-hypotheses} also shows the point null. Denoting both $\Delta_{\text{between}}^{(t)}$ and $\Delta_{\text{within}}^{(t)}$ by $\Delta^{(t)}$, the point null is
    \[
    H_{0}^{(t)}: \Delta^{(t)} = 0 \quad \text{for } t = 1,2,3.
    \]
The unconstrained alternative spans the entire parameter range, which is unbounded and thus covers effect sizes in the range $(-\infty,\infty)$ for all three studies: %(even though the parameters have a different interpretation as between-person vs. within-person differences):
    \[
    H_{u}^{(t)}: \Delta^{(t)} \in (-\infty,\infty) \quad \text{for } t = 1,2,3.
    \]
The complement hypotheses now differ for the three studies, since each study specifies a unique predicted hypothesis. Study 1 specifies the narrowest equivalence interval, and thus its complement starts at a smaller (absolute) value, whereas Study 3 specifies the widest equivalence interval, and thus its complement starts at a larger (absolute) value:
    \[
    \begin{aligned}
    H_{c}^{(1)} &: \Delta^{(1)} \in (-\infty, -0.0625]\ \cup\ [0.0625, \infty) \\
    H_{c}^{(2)} &: \Delta^{(2)} \in (-\infty, -0.1]\ \cup\ [0.1, \infty) \\
    H_{c}^{(3)} &: \Delta^{(3)} \in (-\infty, -0.2]\ \cup\ [0.2, \infty).
    \end{aligned}
    \]


%---------------------------------------
%           Joint Hypotheses           %
%---------------------------------------

\section{Joint hypotheses across studies}

\subsection{Joint predicted hypotheses}
\label{sec:joint-predicted}

The individual studies may test different aspects of the same overarching theory. If the theory holds across different contexts, this implies that the study-specific predicted hypothesis is true in each study. Formally, this joint predicted hypothesis combines all study-specific predicted hypotheses with the logical AND operator:
    \[
    H_1^{(\text{joint})} = H_1^{(1)} \land H_1^{(2)} \land H_1^{(3)} \;\land\; \dots \;\land\; H_1^{(T)}.
    \]
Essentially, the joint hypothesis concatenates the study-specific parameters $\theta^{(t)}$ into a parameter vector
    \[
    \boldsymbol{\theta} = \left(\theta^{(1)}, \, \theta^{(2)}, \, \dots, \, \theta^{(T)}\right),
    \]
where the parameters can be thought of as separate dimensions (indexed per study $t$) of an overarching joint model across all studies.

In BES, the joint predicted hypothesis is tested by multiplying the study-specific Bayes factors, which yields the joint Bayes factor. Previous work has referred to this as the ``product Bayes factor'' \parencite[PBF;][]{vanlissaTutorialAggregatingEvidence2024, klaassenAllOneAll2018}. We use $\mathrm{BF^{(joint)}}$ to avoid conceptual confusion in the upcoming sections:
\begin{equation}
        \mathrm{BF_{1,a}^{(joint)}} = \prod_{t=1}^{T} \mathrm{BF}_{1,a}^{(t)}.
\end{equation}
According to probability theory, simple multiplication is the appropriate operation under three conditions, which must hold for the predicted and the alternative hypotheses alike. First, each study tests a separate hypothesis, meaning that its order and equality constraints concern only its own parameters
$\theta^{(t)}$ and not those of other studies. In this case, the joint hypothesis corresponds to the Cartesian product of the study-specific hypotheses, $H_1^{\text{(joint)}} = H_1^{(1)} \times \dots \times H_1^{(T)}$. Second, the data are conditionally independent across studies given the parameters,
\mbox{$P(D^{(1)}, \ldots, D^{(T)} \mid \boldsymbol{\theta}) = \prod_{t} P(D^{(t)} \mid \theta^{(t)})$}. Third, the prior distributions of the study-specific parameters are independent between studies. Under these assumptions, the joint marginal likelihood factorizes into study-specific marginal likelihoods, and each of these depends only on the hypothesis of its own study:
    \begin{align}
        \mathrm{BF}_{1,a}^{\text{(joint)}}
        &= \frac{P(D \mid H_1^{\text{(joint)}})}{P(D \mid H_a^{\text{(joint)}})}
        \nonumber \\
        &= \frac
          {P(D^{(1)}, \dots, D^{(T)} \mid H_1^{(1)} \land \dots \land H_1^{(T)})}
          {P(D^{(1)}, \dots, D^{(T)} \mid H_a^{(1)} \land \dots \land H_a^{(T)})}
        \nonumber \\
        &= \prod_{t=1}^{T}
          \frac{P(D^{(t)} \mid H_1^{(1)} \land \dots \land H_1^{(T)})}
               {P(D^{(t)} \mid H_a^{(1)} \land \dots \land H_a^{(T)})}
        \nonumber \\
        &= \prod_{t=1}^{T} \frac{P(D^{(t)} \mid H_1^{(t)})}{P(D^{(t)} \mid H_a^{(t)})}
        \label{eq:bf_1a_joint}.
    \end{align}

\subsubsection{Example A: Choice probabilities}

In Example A, the joint predicted hypothesis states that the hit rate for positive stimuli is larger than 0.50 in Study 1 \textit{and} Study 2 \textit{and} Study 3:
    \[
    H_{1}^{(\text{joint})}: (\theta^{(1)} > 0.50) \;\land\; (\theta^{(2)} > 0.50) \;\land\; (\theta^{(3)} > 0.50).
    \]
This joint hypothesis can be visualized in a multidimensional parameter space, in which the conjunction of multiple order-constrained hypotheses creates a restricted subspace \parencite[i.e., a convex polytope;][]{heckMultinomialModelsLinear2019}. For Example A, this subspace corresponds to the pink cube in each of the three panels in the left column of Figure~\ref{fig:joint-hypotheses}; the other shaded volumes represent different joint alternative hypotheses, which we discuss in the next sections.

\begin{figure}[!tbp]
\centering

\captionsetup{justification=raggedright, singlelinecheck=false}
\caption{Joint hypotheses for Example A (left column) and Example B (right column)}
\label{fig:joint-hypotheses}
\vspace{0pt}

\begin{tabular}{c @{\hspace{30pt}} !{\vrule width 0.5pt} @{\hspace{30pt}} c}

\begin{subfigure}{0.40\textwidth}
    \centering
    \caption*{(a)}
    \includegraphics[width=\linewidth]{figures/fig2_A2_choice_prob_hu.pdf}
\end{subfigure}
&
\begin{subfigure}{0.40\textwidth}
    \centering
    \caption*{(d)}
    \includegraphics[width=\linewidth]{figures/fig2_B2_equivalence_hu.pdf}
\end{subfigure}
\\
\begin{subfigure}{0.40\textwidth}
    \centering
    \caption*{(b)}
    \includegraphics[width=\linewidth]{figures/fig2_A3_choice_prob_incomplete_hc.pdf}
\end{subfigure}
&
\begin{subfigure}{0.40\textwidth}
    \centering
    \caption*{(e)}
    \includegraphics[width=\linewidth]{figures/fig2_B3_equivalence_incomplete_hc.pdf}
\end{subfigure}
\\
\begin{subfigure}{0.40\textwidth}
    \centering
    \caption*{(c)}
    \includegraphics[width=\linewidth]{figures/fig2_A4_choice_prob_complete_hc.pdf}
\end{subfigure}
&
\begin{subfigure}{0.40\textwidth}
    \centering
    \caption*{(f)}
    \includegraphics[width=\linewidth]{figures/fig2_B4_equivalence_complete_hc.pdf}
\end{subfigure}

\end{tabular}

\vspace{15pt}

\begin{minipage}{\textwidth}
\raggedright
\footnotesize
\textit{Note.} Joint predicted hypothesis (pink cuboid) with the joint null (blue point, all panels) and one further joint alternative per panel: unconstrained (yellow; a, d), incomplete complement (green; b, e), and complete complement (green; c, f). For Example~B, the unbounded parameter space is truncated.
\end{minipage}

\end{figure}

\subsubsection{Example B: Equivalence testing}

In Example B, the joint predicted hypothesis states that, in each study, the mean difference in test scores between the identity-primed and the non-primed condition falls within its respective equivalence interval:
    \[
    H_{1}^{(\text{joint})}:
    \left(-0.0625 < \Delta^{(1)} < 0.0625\right)
    \;\land\;
    \left(-0.1 < \Delta^{(2)} < 0.1\right)
    \;\land\;
    \left(-0.2 < \Delta^{(3)} < 0.2\right).
    \]
The joint predicted hypothesis now forms a cuboid rather than a cube as in Example A (right column of Figure~\ref{fig:joint-hypotheses}), because the equivalence intervals of the studies span different proportions of their parameter ranges.

\subsection{Joint alternative hypotheses}
\label{sec:joint-alternative-hypotheses}

Evidence for the joint predicted hypothesis is accumulated in the numerator of $\text{BF}_{1,a}^{(\text{joint})}$ in Eq.~\eqref{eq:bf_1a_joint}, which reflects how plausible the observed data of all studies together are if the predictions hold in all studies. This raises the question of what the denominator of the joint Bayes factor represents.

The joint alternative hypothesis states that the study-specific alternative hypothesis is true in each study:
    \[
    H_a^{(\text{joint})} = H_a^{(1)} \;\land\; H_a^{(2)} \;\land\; H_a^{(3)} \;\land\; \dots \;\land\; H_a^{(T)}.
    \]
This joint alternative arises automatically when multiplying the study-specific Bayes factors. Although this follows directly from the formula for the product Bayes factor, the meaning and interpretation of different joint alternative hypotheses have, to the best of our knowledge, not previously been compared and discussed in detail. We consider three cases: all studies test against the null hypothesis, all studies test against the unconstrained hypothesis, and all studies test against the complement hypothesis.

\subsubsection{Example A: Choice probabilities}

If all three replication studies test $H^{(t)}_1$ against the null, $H^{(t)}_0: \theta^{(t)} = 0.50$, then the joint alternative hypothesis states that the hit rate is 50\% in all studies:
    \[
    H_{0}^{(\text{joint})}: \boldsymbol{\theta} = (0.50, 0.50, 0.50).
    \]
The joint null hypothesis corresponds to a single point in the multidimensional parameter space (the blue point in each of the three plots in the left column of Figure~\ref{fig:joint-hypotheses}). If all three studies test against the unconstrained alternative, $H_{u}^{(t)}: \theta^{(t)} \in [0,1]$, then the joint alternative hypothesis simply includes all possible constellations of $\theta^{(t)}$ across studies, that is, in each study, the parameters can take on any value in the range $[0,1]$:
    \[
    H_{u}^{(\text{joint})}: \boldsymbol{\theta} = \left(\theta^{(1)}, \theta^{(2)}, \theta^{(3)}\right) \in [0,1]^3.
    \]
The joint unconstrained hypothesis thus corresponds to the entire joint parameter space (yellow volume in Figure~\ref{fig:joint-hypotheses}a). Finally, if all three studies test against the complement alternative $H^{(t)}_c: \theta^{(t)} \leq 0.50$, then the joint alternative hypothesis states that $\theta^{(t)} \leq 0.50$ in all studies:
    \[
    H_{c}^{(\text{joint})}: \left(\theta^{(1)} \leq 0.50\right) \;\land\; \left(\theta^{(2)} \leq 0.50\right) \;\land\; \left(\theta^{(3)} \leq 0.50\right),
    \]
or equivalently, $\boldsymbol{\theta} \in [0, 0.50]^3$ (green volume in Figure~\ref{fig:joint-hypotheses}b).

\subsubsection{Example B: Equivalence testing}

In Example B, the joint null states that $\Delta^{(t)}=0$ in each study:
    \[
    H_{0}^{(\text{joint})}: \boldsymbol{\Delta} = (0,0,0).
    \]
As in Example~A, this corresponds to a single point in the joint parameter space (the blue point in each of the three plots in the right column of Figure~\ref{fig:joint-hypotheses}). The joint unconstrained hypothesis again covers the entire parameter space (yellow volume in Figure~\ref{fig:joint-hypotheses}d):
    \[
    H_{u}^{(\text{joint})}: \boldsymbol{\Delta} =
    (\Delta^{(1)}, \Delta^{(2)}, \Delta^{(3)}) \in \mathbb{R}^3.
    \]
The joint complement is the conjunction of all study-specific complement hypotheses (green volume in Figure~\ref{fig:joint-hypotheses}e):
    \[
        H_{c}^{(\text{joint})}:
        \left(\left|\Delta^{(1)}\right| \geq 0.0625\right)
        \land
        \left(\left|\Delta^{(2)}\right| \geq 0.1\right)
        \land
        \left(\left|\Delta^{(3)}\right| \geq 0.2\right).
    \]
% In contrast to Example A, this region no longer forms a single connected volume. It is instead the union of eight disjoint cuboids $\Theta_i$, three of which are
%     \begin{align*}
%         \Theta_1 &= (-\infty, -0.0625] \times (-\infty, -0.1] \times (-\infty, -0.2]\\
%         \dots \\
%         \Theta_3 &= [0.0625, \infty) \times [0.1, \infty) \times (-\infty, -0.2] \\
%         \dots \\
%         \Theta_8 &= (-\infty, -0.0625] \times [0.1, \infty) \times [0.2, \infty).
%    \end{align*}



%----------------------------------------------------------------
%           Interpretation of the Joint Bayes Factor          %
%----------------------------------------------------------------

\section{Interpretation of the joint Bayes factor}

\label{sec:interpretation-joint-bf}

The $\mathrm{BF}_{1,a}^{\text{(joint)}}$ pits the joint predicted against a joint alternative hypothesis. The correct interpretation of the joint Bayes factor depends on the choice of alternative hypothesis. If each study tests the predicted against a point null hypothesis, the resulting joint Bayes factor $\mathrm{BF}_{1,0}^{\text{(joint)}}$ can be interpreted as the relative evidence that the theory holds in each individual study compared to the effect being absent in each study. If we test against the unconstrained hypothesis, then the product Bayes factor $\mathrm{BF}_{1,u}^{\text{(joint)}}$ pits the joint $H_1$ against all possibilities of parameter combinations across studies. The joint Bayes factor can then be interpreted as relative evidence that the theory holds in each individual study compared to \textit{not having any theoretical account for the data}. Framed differently, the Bayes factor quantifies the degree to which imposing theoretical constraints improves our predictions of the data compared to not making any predictions at all.

Finally, if each study tests the predicted hypothesis against its complement, the resulting joint Bayes factor $\mathrm{BF}_{1,c}^{\text{(joint)}}$ pits $H_1^{\text{(joint)}}$ against the joint complement $H_c^{\text{(joint)}}$ formed by combining the study-specific complements. In this third case, however, the joint Bayes factor is difficult to interpret. A tempting misinterpretation is that $\mathrm{BF}_{1,c}^{\text{(joint)}}$ quantifies the relative evidence for the set of theoretical predictions versus \textit{its} complement. This interpretation is incorrect because the conjunction of the study-specific complements does not form the complement of $H_1^{\text{(joint)}}$ in the multidimensional parameter space. The mismatch is immediately visible in Panels~b and~e of Figure~\ref{fig:joint-hypotheses}: Outside the pink cuboids representing $H_1^{\text{(joint)}}$, large regions of the parameter space are not covered by the green cuboids representing $H_c^{\text{(joint)}}$. The joint alternative that arises from multiplying the study-specific Bayes factors is therefore an \textit{incomplete} complement of the theory, which does not represent a theoretically meaningful claim. Correctly interpreted, $\mathrm{BF}_{1,c}^{\text{(joint)}}$ quantifies the relative evidence that the theoretical predictions hold in every study versus that they are false in every study.

The complete complement of $H_1^{\text{(joint)}}$, by contrast, covers all combinations of parameter values \textit{except} those covered by $H_1^{\text{(joint)}}$. It therefore includes the incomplete complement as the extreme case in which every study-specific complement is true, but also many other combinations, such as $H_1^{(t)}$ holding in only one or two studies and $H_c^{(t)}$ in the remaining ones. In other words, the complete complement states that the theoretical predictions are false in at least one study. Geometrically, it corresponds to the entire joint parameter space minus the pink region, as shown in Panels~c and~f of Figure~\ref{fig:joint-hypotheses}.

%------------------------------------------------------------------------
%           Computing the Evidence for the Complete Complement          %
%------------------------------------------------------------------------

\section{Computing the evidence for the complete complement}
\label{sec:compute-complete-complement}

In order for the joint Bayes factor to quantify the relative evidence for the theory against \textit{its complement}, the complete joint complement $H_{c^{*}}^{\text{(joint)}}$ must be considered as the alternative. In this section, we show how the resulting Bayes factor $\mathrm{BF}_{1,c^{*}}^{\text{(joint)}}$ can be computed from the study-specific Bayes factors, that is, from the same information required for standard BES.%
    \footnote{Strictly speaking, the superscript ``(joint)'' is redundant here because the distinction between the complete and the incomplete complement arises only in the joint parameter space: At the study level, $H_c^{(t)}$ is already the complete complement of $H_1^{(t)}$. We keep the superscript for consistency with the other joint hypotheses.}

The computations of $\mathrm{BF}_{1,c^{*}}^{\text{(joint)}}$ are based on the encompassing prior approach for order-constrained hypotheses \parencite{klugkistBayesianModelSelection2005, klugkistBayesFactorInequality2007}, in which each informative hypothesis $H_i$ is a nested model within the unconstrained hypothesis. Once a prior distribution is specified across the full parameter space of the unconstrained model (the so-called encompassing prior), $\mathrm{BF}_{i,u}$ can be computed as the ratio of two quantities: the fit $f_i$, which is the proportion of the posterior distribution in agreement with the order constraints defined by $H_i$, and the complexity $c_i$, which is the corresponding proportion of the prior distribution. Since the complexity of an equality-constrained hypothesis is zero under a continuous prior, this ratio is undefined for such hypotheses, so the following computations apply only to predicted hypotheses defined by order and interval constraints. For each study $t$, this yields
\begin{align}
\mathrm{BF}_{1,u}^{(t)}
&= \frac{P\left(D^{(t)} \mid H_1^{(t)}\right)}
        {P\left(D^{(t)} \mid H_u^{(t)}\right)}\nonumber \\
&= \frac{P\left(\theta^{(t)} \in H_1^{(t)} \mid D^{(t)}\right)}
        {P\left(\theta^{(t)} \in H_1^{(t)}\right)}
= \frac{f_1^{(t)}}{c_1^{(t)}}. 
\label{eq:bf-1u-t}
\end{align}

To illustrate the following steps concretely, we return to Example A. Suppose the three replications of Bem's Experiment 1 yielded the study-specific Bayes factors $\mathrm{BF}_{1,u}^{(1)} = 1.80$, $\mathrm{BF}_{1,u}^{(2)} = 1.50$, and
$\mathrm{BF}_{1,u}^{(3)} = 0.70$.%
    % \footnote{These values are hypothetical but
    % plausible: since $H_1^{(t)}: \theta^{(t)} > 0.50$ has complexity
    % $c_1^{(t)} = 0.50$, the evidence in favor of the prediction is bounded above by
    % $1/c_1^{(t)} = 2$.}
The first two studies thus provide weak evidence in favor of precognition for positive stimuli, whereas the third slightly favors the
unconstrained alternative.

\subsection{Step 1: Compute the Bayes factor for the complement in each study}
\label{sec:step-1-compute-bf-c-u-t}

Assume that each study provides only the Bayes factor $\mathrm{BF}_{1,u}^{(t)}$. First, we obtain, for each study, the Bayes factor $\mathrm{BF}_{c,u}^{(t)}$ of the complement against the unconstrained hypothesis. Sometimes, $\mathrm{BF}_{c,u}^{(t)}$ is directly available in software output, where the study-specific complement is considered automatically (e.g., the R packages \texttt{bain} by \cite{guBainBayesFactors2024}, \texttt{BFpack} by \cite{mulderBFpackFlexibleBayes2021}, and \texttt{multinomineq} by \cite{heckMultinomineqBayesianInference2019}). To manually compute this Bayes factor, it is necessary to derive or compute the complexity $c_1^{(t)}$ of the predicted hypothesis $H_1^{(t)}$, that is, the prior probability $P\left(\theta^{(t)} \in H_1^{(t)}\right)$ of the parameters being in the predicted hypothesis space. Applying the encompassing Bayes factor in Eq.~\eqref{eq:bf-1u-t} for the complement gives
\begin{align}
\mathrm{BF}_{c,u}^{(t)}
&= \dfrac{P\left(D^{(t)} \mid H_c^{(t)}\right)}
         {P\left(D^{(t)} \mid H_u^{(t)}\right)}
 = \dfrac{1-f_1^{(t)}}{1-c_1^{(t)}} \nonumber \\
&= \dfrac{1-P\left(\theta^{(t)} \in H_1^{(t)} \mid D^{(t)}\right)}
         {1-P\left(\theta^{(t)} \in H_1^{(t)}\right)} \nonumber \\
&= \dfrac{1-c_1^{(t)}\,\mathrm{BF}_{1,u}^{(t)}}{1-c_1^{(t)}}.
\label{eq:bf-cu-t}
\end{align}

Returning to Example A, and assuming a uniform prior on $\theta^{(t)}$ (or any other prior that is symmetric around 0.50), the complexity of $H_1^{(t)}: \theta^{(t)} > 0.50$ is $c_1^{(t)} = 0.50$ in all three studies, since half of the prior distribution satisfies the order constraint. Applying Eq.~\eqref{eq:bf-cu-t} to the Bayes factor of Study~1, $\mathrm{BF}_{1,u}^{(1)} = 1.80$, gives
\[
\mathrm{BF}_{c,u}^{(1)} = \frac{1 - 0.50 \cdot 1.80}{0.50} = 0.20.% \quad
%\mathrm{BF}_{c,u}^{(2)} = \frac{1 - 0.50 \times 1.50}{0.50} = 0.50, \quad
%\mathrm{BF}_{c,u}^{(3)} = \frac{1 - 0.50 \times 0.70}{0.50} = 1.30.
\]
Similarly, we obtain $\mathrm{BF}_{c,u}^{(2)} =0.50$ and $\mathrm{BF}_{c,u}^{(3)} =1.30$.

\subsection{Step 2: Compute the Bayes factor for each part of the complete complement}

\label{sec:step-2-compute-bf-omega-u-t}

Knowing both $\mathrm{BF}_{1,u}$ and $\mathrm{BF}_{c,u}$ allows us to construct all possible combinations of study-specific predicted and complement hypotheses across studies. That is, instead of only considering those two parts of the joint parameter space where either $H_1$ is true in \textit{all} studies or $H_c$ is true in \textit{all} studies, we now also consider all combinations where $H_1$ holds in \textit{at least one, but not all} studies.

To achieve this, we enumerate each possible combination by defining the index set $\Omega \subseteq \{1,\dots, T\}$ as the set of studies in which the complement $H_c^{(t)}$ is assumed to hold. In the remaining studies $t \notin \Omega$, the predicted hypothesis $H_1^{(t)}$ is assumed to hold. For instance, $\Omega = \{1,3\}$ denotes the case in which the complement holds in Studies~1 and~3 while the theory-derived predictions are correct in all other studies. Since each study can either be included in or excluded from $\Omega$, the number of possible unique sets $\Omega$ is
$$
\left|\left\{\Omega : \Omega \subseteq \{1,\dots,T\}\right\}\right| = 2^T.
$$
Each set $\Omega$ defines a specific joint hypothesis $H_{\Omega}$ that corresponds to a unique non-overlapping region in the multidimensional parameter space. 
%Except $\Omega = \emptyset$ (i.e., the predictions are true in each study), each set constitutes one component of the complete complement. 

We now compute a product Bayes factor that quantifies the relative evidence provided by all studies together in favor of the specific hypothesis $H_\Omega$ relative to the unconstrained hypothesis:
\begin{align}
\mathrm{BF}_{\Omega,u}^{\text{(joint)}}
&= \frac{P\left(D \mid H_{\Omega}\right)}
        {P\left(D \mid H_u\right)} \nonumber \\
&= \frac{\prod_{t} P\left(D^{(t)} \mid H_{\Omega}\right)}
        {\prod_{t} P\left(D^{(t)} \mid H_u^{(t)}\right)} \nonumber \\
&= \frac{\prod_{t \in \Omega} P\left(D^{(t)} \mid H_c^{(t)}\right)
         \prod_{t \notin \Omega} P\left(D^{(t)} \mid H_1^{(t)}\right)}
        {\prod_{t} P\left(D^{(t)} \mid H_u^{(t)}\right)} \nonumber \\
&= \prod_{t \in \Omega} \mathrm{BF}_{c,u}^{(t)} \cdot \prod_{t \notin \Omega} \mathrm{BF}_{1,u}^{(t)}.
\label{eq:bf-omega-u-joint}
\end{align}
Note that this formulation includes two special cases already discussed above. If $\Omega = \{1,\dots,T\}$, the complement $H_c$ is assumed to hold in all studies. Using the convention that the product over an empty set equals 1, we obtain the joint Bayes factor that pits the incomplete joint complement (i.e., the conjunction of study complements) against the joint unconstrained hypothesis:
$$
\mathrm{BF}_{c,u}^{\text{(joint)}} = \prod_{t \in \{1,\dots,T\}} \mathrm{BF}_{c,u}^{(t)}.
$$
The second case, $\Omega = \emptyset$, reflects the assumption that $H_1$ holds in all studies, which yields the joint Bayes factor that pits the joint predicted against the joint unconstrained hypothesis:
$$
\mathrm{BF}_{1,u}^{\text{(joint)}} = \prod_{t \in \{1,\dots,T\}} \mathrm{BF}_{1,u}^{(t)}.
$$

In Example~A, with $T = 3$ studies, there are $2^T = 8$ index sets $\Omega$, each corresponding to one of the eight small cubes in Figure~\ref{fig:joint-hypotheses}c. 
For instance, $\Omega = \{2,3\}$ corresponds to the lower right cube, reflecting the hypothesis $H_{{\Omega}}$ that the parameter $\theta$ is above 0.50 (as predicted by precognition) only in Study~1, but below 0.50 in Studies~2 and~3.
Table~\ref{tab:example_a_omega} lists the resulting Bayes factors $\mathrm{BF}_{\Omega,u}^{\text{(joint)}}$ computed via Eq.~\eqref{eq:bf-omega-u-joint}.

\begin{table}[!tb]
\captionsetup{justification=raggedright, singlelinecheck=false}
\caption{Hypothetical Bayes factors for all combinations of $H_1$ and $H_c$ in Example A}
\label{tab:example_a_omega}
\centering
\begin{tabular}{l>{\centering\arraybackslash}c>{\centering\arraybackslash}c>{\centering\arraybackslash}c}
\hline
$\Omega$ & $H_\Omega$ & Computation & $\mathrm{BF}_{\Omega,u}^{\text{(joint)}}$ \\
\hline
$\emptyset$   & $H_1 \wedge H_1 \wedge H_1$ & $1.80 \times 1.50 \times 0.70$ & 1.89 \\
$\{1\}$       & $H_c \wedge H_1 \wedge H_1$ & $0.20 \times 1.50 \times 0.70$ & 0.21 \\
$\{2\}$       & $H_1 \wedge H_c \wedge H_1$ & $1.80 \times 0.50 \times 0.70$ & 0.63 \\
$\{3\}$       & $H_1 \wedge H_1 \wedge H_c$ & $1.80 \times 1.50 \times 1.30$ & 3.51 \\
$\{1,2\}$     & $H_c \wedge H_c \wedge H_1$ & $0.20 \times 0.50 \times 0.70$ & 0.07 \\
$\{1,3\}$     & $H_c \wedge H_1 \wedge H_c$ & $0.20 \times 1.50 \times 1.30$ & 0.39 \\
$\{2,3\}$     & $H_1 \wedge H_c \wedge H_c$ & $1.80 \times 0.50 \times 1.30$ & 1.17 \\
$\{1,2,3\}$   & $H_c \wedge H_c \wedge H_c$ & $0.20 \times 0.50 \times 1.30$ & 0.13 \\
\hline
\end{tabular}
\vspace{2mm}
\begin{flushleft}
\footnotesize
\textit{Note.} $\Omega$ is the set of studies in which $H_c^{(t)}$ is assumed to hold, whereas $H_1^{(t)}$ is assumed to hold in the remaining studies. Because $c_1^{(t)} = 0.50$ in all studies, Eq.~\eqref{eq:bf-cu-t} implies that $\mathrm{BF}_{1,u}^{(t)}$ and $\mathrm{BF}_{c,u}^{(t)}$ cannot exceed 2 and sum to 2.
\end{flushleft}
\end{table}

\subsection{Step 3: Compute the Bayes factor for the complete complement}
\label{sec:step-3-compute-bf-cstar-u-joint}

The complete joint complement is defined by all non-empty sets $\Omega$, meaning that the theoretical predictions are incorrect in at least one study. This is best seen in Figure~\ref{fig:joint-hypotheses}c, where the complement consists of seven green cubes for $T=3$ studies. The Bayes factor of the complete complement against the joint unconstrained hypothesis is obtained as a weighted average of the Bayes factors
$\mathrm{BF}^{(\text{joint})}_{\Omega,u}$ across all non-empty sets
$\Omega$:
    \footnote{This also explains why we avoid the term ``PBF.'' Unlike the other joint Bayes factors, $\mathrm{BF}^{(\text{joint})}_{1,c^*}$ is not a product of study-specific Bayes factors.}
\begin{equation}
\label{eq:bf-cstar-u-joint}
\mathrm{BF}^{(\text{joint})}_{c^*,u} =
\sum_{\emptyset\neq\Omega\subseteq\{1,\ldots,T\}} \pi_\Omega
\prod_{t\in\Omega}\mathrm{BF}^{(t)}_{c,u}\prod_{t\notin\Omega}\mathrm{BF}^{(t)}_{1,u}.
\end{equation}

The weights $\pi_\Omega$ are the prior probabilities of the individual components $H_\Omega$ under the assumption that the complete complement holds. These
weights express how plausible each pattern of successful and failed predictions across studies is
considered a priori. There are two natural choices for the weights.
\textit{Equal weighting} assigns the same weight to every component of the complete complement,
  \begin{equation}
  \pi_\Omega = \frac{1}{2^T-1},
  \end{equation}
so that all patterns of successful and failed predictions are considered equally plausible a priori.
In contrast, \textit{complexity weighting} assigns each component a weight equal to its share of the prior mass of the complete complement,
  \begin{equation}
  \pi_\Omega
  = \frac{c_\Omega}{\sum_{\emptyset\neq\Omega'\subseteq\{1,\dots,T\}} c_{\Omega'}}
  = \frac{c_\Omega}{1-\prod_{t=1}^{T} c_1^{(t)}},
  \qquad
  c_\Omega = \prod_{t\in\Omega}\bigl(1-c_1^{(t)}\bigr)\prod_{t\notin\Omega} c_1^{(t)},
  \end{equation}
where $c_\Omega$ is the complexity of the component $H_\Omega$ and $c_1^{(t)}$ the complexity of the predicted hypothesis in study~$t$. This corresponds to restricting the encompassing prior to the complete complement.

By the transitivity of Bayes factors in Eq.~\eqref{eq:transitivity}, $\mathrm{BF}^{(\text{joint})}_{1,c^*} = \mathrm{BF}^{(\text{joint})}_{1,u}/\mathrm{BF}^{(\text{joint})}_{c^*,u}$. Substituting the expressions for the weights into Eq.~\eqref{eq:bf-cstar-u-joint} yields simplified expressions for $\mathrm{BF}^{(\text{joint})}_{1,c^*}$ that can be computed directly from the study-specific quantities, without summing over all $2^T-1$ components (for detailed derivations, see Appendix~\ref{app:closed-form}). Under complexity weighting, $\mathrm{BF}^{(\text{joint})}_{1,c^*}$ depends only
on $\mathrm{BF}^{(t)}_{1,u}$ and $c_1^{(t)}$:
\begin{equation}
\label{eq:bf-1-cstar-complexity}
\mathrm{BF}^{(\text{joint})}_{1,c^*}
= \frac{\left(1-\prod_{t=1}^{T} c_1^{(t)}\right)\prod_{t=1}^{T}\mathrm{BF}^{(t)}_{1,u}}
       {1-\prod_{t=1}^{T}\left(c_1^{(t)}\,\mathrm{BF}^{(t)}_{1,u}\right)}.
\end{equation}
Under equal weighting, $\mathrm{BF}^{(\text{joint})}_{1,c^*}$ depends only on $\mathrm{BF}^{(t)}_{1,c}$:
\begin{equation}
\label{eq:bf-1-cstar-equal}
\mathrm{BF}^{(\text{joint})}_{1,c^*}
= \frac{2^T-1}{\prod_{t=1}^{T}\left(1 + 1/\mathrm{BF}^{(t)}_{1,c}\right) - 1}.
\end{equation}
The study-specific Bayes factors used in Eq.~\eqref{eq:bf-1-cstar-equal} are exactly the ones that are multiplied in BES when testing against the incomplete complement. The two types of weighting yield identical results whenever $c_1^{(t)} = 0.50$ in all studies, because all components then have the same prior mass $c_\Omega = 2^{-T}$.

\subsection{Choosing prior weights for the complete complement}

Equal and complexity weighting differ in how the prior mass of the complete complement is distributed across its components. We use equal weighting throughout, because it accords with the premise of BES that each study provides an equally valid test of the same theory. By Eq.~\eqref{eq:bf-1-cstar-equal}, each study affects $\mathrm{BF}^{(\text{joint})}_{1,c^*}$ only through the Bayes factor $\mathrm{BF}^{(t)}_{1,c}$ for the predicted hypothesis against its complement. Therefore, equally strong evidence against the prediction is equally detrimental to the theory, regardless of which study provides it. In contrast, under complexity weighting, the complete complement expects failed predictions primarily in those studies with predicted hypotheses of low complexity (i.e., hypotheses that cover only a small part of the prior), so that evidence against these predictions counts more strongly. For instance, consider two studies whose predicted hypotheses have complexities $c_1 = 0.10$ and $c_1 = 0.50$, and assume that one of them yields $\mathrm{BF}_{1,c} = 1/3$, whereas the other provides no evidence ($\mathrm{BF}_{1,c} = 1$). Under equal weighting, $\mathrm{BF}^{(\text{joint})}_{1,c^*} = 3/7 \approx 0.43$, regardless of which study provides the evidence. Under complexity weighting, the joint Bayes factor is $0.35$ if the evidence comes from the study with $c_1 = 0.10$ but $0.49$ if it comes from the study with $c_1 = 0.50$.

The asymmetric impact of studies on $\mathrm{BF}^{(\text{joint})}_{1,c^*}$ that comes with complexity weighting is defensible if all studies concern the same parameter on the same scale, because a narrower predicted hypothesis then represents a more precise, and hence riskier, prediction. In BES, however, studies often concern different parameters with different ranges, so that differences in complexity need not reflect differences in how precisely the theory is tested. For instance, in Example~B, the equivalence intervals differ only because one grade point corresponds to a different number of percentage points in each study's grading system. With the same prior on the mean difference in all three studies, Study 1 would have the smallest complexity, so complexity weighting would treat it as the most likely site of falsification and let evidence against its prediction count most strongly. Yet there is no substantive reason why a failed prediction in Study 1 should weigh more heavily against the theory than one in Study 3. Equal weighting avoids such asymmetries.

\subsection{Joint complement hypothesis for multiple predicted hypotheses}

So far, we have assumed that each study tests a single predicted hypothesis. Steps~1 to~3 also apply to scenarios in which each study under consideration tests multiple competing predicted hypotheses (e.g., qualitatively different sets of order constraints on the parameters). As in the single-hypothesis case, multiplying the study-specific Bayes factors for one of the predicted hypotheses against its complement does not yield a test against the complete complement. Unlike in that case, however, considering a separate study-specific complement hypothesis for each predicted hypothesis would result in a very large set of hypotheses, many of which overlap substantially in the parameter space. 

To compare multiple joint predicted hypotheses simultaneously (e.g., via posterior model probabilities), it is more sensible to define a single complement hypothesis that includes all regions of the parameter space not covered by \textit{any} of the joint predicted hypotheses. The complete complement across studies then comprises all combinations of study-specific hypotheses in which none of the predicted hypotheses holds across all studies. Section~S2 of the Supplementary Material (\suppurl) derives the corresponding Bayes factors and posterior model probabilities.

% Once we know the BF for each of two constrained hypotheses against the unconstrained
% hypothesis, the BF comparing these two directly follows from transitivity. Therefore, the BF comparing the joint predicted hypothesis $H_1^{(\text{joint})}$ against its complement
% $H_{c^*}^{(\text{joint})}$ can be computed as
% \begin{equation}
% \label{eq:bf-1-cstar}
% \mathrm{BF}^{(\text{joint})}_{1,c^*} =
% \frac{\mathrm{BF}^{(\text{joint})}_{1,u}}{\mathrm{BF}^{(\text{joint})}_{c^*,u}}.
% \end{equation}
% Applying Eq.~\eqref{eq:bf-cstar-u-joint} with equal weights to Example A, the complete
% joint complement is the average of the seven rows of non-empty sets $\Omega$ in
% Table~\ref{tab:exampleA},
% \begin{equation*}
% \mathrm{BF}^{(\text{joint})}_{c^*,u} =
% \frac{1}{7}\,(0.21 + 0.63 + 3.51 + 0.07 + 0.39 + 1.17 + 0.13) = 0.87,
% \end{equation*}
% and Eq.~\eqref{eq:bf-1-cstar} yields
% $\mathrm{BF}^{(\text{joint})}_{1,c^*} = 1.89/0.87 = 2.17$. Because $c_1^{(t)} = 0.5$ in
% all three studies, complexity weighting yields the same result.


%-------------------------------------
%           Empirical Examples       %
%-------------------------------------
\section{Empirical examples}

In this section, we investigate the extent to which different joint alternative hypotheses lead to different conclusions about the theory of interest in two empirical examples. In both cases, we computed $\mathrm{BF}_{1,c^{*}}^{(\text{joint})}$ with equal weighting from the study-specific $\mathrm{BF}_{1,u}^{(t)}$, following Steps~1--3 in Section~\ref{sec:compute-complete-complement}. Data and R scripts for the analyses are available at \url{https://github.com/L-Groot/BES}.

\subsection{Hotel towel reuse% \texorpdfstring{\parencite{scheibehenneBayesianEvidenceSynthesis2016}}{(Scheibehenne et al., 2016)}
}

\textcite{scheibehenneBayesianEvidenceSynthesis2016} investigated whether descriptive social norms increase pro-environmental behavior. They considered four published papers (some with multiple experiments), which all tested whether normative messages encourage hotel guests to reuse their towels. Each experiment tested the prediction that guests exposed to a descriptive social-norm message are more likely to reuse their towels than guests in the control condition; that is, the proportion of guests reusing their towels is expected to be higher in the experimental than in the control condition ($H_1: \theta_e > \theta_c$). In the majority of experiments (five out of seven), the Bayes factor $\mathrm{BF}_{1,0}^{(t)}$ favored the null hypothesis that the two proportions are equal ($H_0: \theta_e = \theta_c$), some even with strong evidence. However, when the data from all seven experiments were pooled and analyzed jointly, the resulting $\mathrm{BF}_{1,0}=36.89$ provided strong evidence \textit{against} the null, in favor of the social-norm hypothesis.

The pooled analysis assumes a fixed-effect model, which implies that all studies share the same effect size. However, there is considerable heterogeneity across the included studies. For example, in some experiments, the control message contained an appeal to the environmental benefit of towel reuse \parencite{goldsteinRoomViewpointUsing2008, bohnerRoomViewpointRevisited2014}, whereas in others it simply instructed guests on where to leave their towels if they wished them to be replaced \parencite{schultzUsingNormativeSocial2008}. Furthermore, one study presented the descriptive norm together with an injunctive norm \parencite{schultzUsingNormativeSocial2008}. In light of this heterogeneity, one may decide against assuming a fixed effect across all studies and apply BES to synthesize the evidence instead. 

    \begin{table}[!tb]
    \captionsetup{justification=raggedright, singlelinecheck=false}
    \caption{Study-specific and joint Bayes factors testing whether normative messages encourage hotel guests to reuse their towels \parencite{scheibehenneBayesianEvidenceSynthesis2016}}
    \label{tab:scheibehenne}
    \centering
    \begin{tabular}{l>{\centering\arraybackslash}c>{\centering\arraybackslash}c>{\centering\arraybackslash}c}
    \hline
    Study &
    $\mathrm{BF}_{1,0}$ &
    $\mathrm{BF}_{1,u}$ &
    $\mathrm{BF}_{1,c}$ \\
    \hline
    Goldstein, Cialdini, \& Griskevicius (2008), Exp. 1 & 1.45 & 1.94 & 32.67\\
    Goldstein et al. (2008), Exp. 2 & 2.02 & 1.97 & 75.34\\
    Schultz, Khazian, \& Zaleski (2008), Exp. 2 & 0.35 & 1.72 & 6.15 \\
    Schultz et al. (2008), Exp. 3 & 0.58 & 1.86 & 13.04\\
    Mair \& Bergin-Seers (2010) & 0.32 & 1.22 & 1.58\\
    Bohner \& Schlüter (2014), Exp. 1 & 0.06 & 0.67 & 0.50\\
    Bohner \& Schlüter (2014), Exp. 2 & 0.06 & 0.04 & 0.02\\
    \hline
    BES: product BF
    & $7.94 \times 10^{-4}$
    & $0.36$
    & $2.85 \times 10^{3 \, \dagger}$ \\
    BES: complete complement
    &
    &
    & $0.36^{*}$ \\
    \hline
    \end{tabular}
    \vspace{2mm}
    \begin{flushleft}
    \footnotesize
    \textit{Note.} Study-specific hypotheses: $H_1: \theta_e > \theta_c$, $H_0: \theta_e = \theta_c$, $H_u: \theta_e, \theta_c \in [0,1]$, and $H_c: \theta_e \leq \theta_c$. Following \textcite[p.~1043]{scheibehenneBayesianEvidenceSynthesis2016}, $\theta_e$ and $\theta_c$ have independent uniform priors under $H_u$; hence, $c_1 = 0.50$ in all experiments. $^{\dagger}$Evidence against the incomplete joint complement. $^{*}$Evidence against the complete joint complement. Some $\mathrm{BF}_{1,0}$ differ by 0.01 from those reported by \textcite{scheibehenneBayesianEvidenceSynthesis2016} due to random sampling from the posterior.
    \end{flushleft}
    \end{table}

Table~\ref{tab:scheibehenne} shows the study-specific and joint Bayes factors for each of the four joint alternative hypotheses discussed in Sections~\ref{sec:joint-alternative-hypotheses} and~\ref{sec:compute-complete-complement}. Only the comparison to the incomplete joint complement favors the joint predicted hypothesis, with $\mathrm{BF}_{1,c}^{\text{(joint)}}=2.85 \times 10^3$ indicating overwhelming support. Thus, the data are much more likely if the theory-derived predictions hold in each study than if they are wrong in each study. However, if the substantive question is whether the overall theory holds across \emph{all} studies or is incorrect in at least one study, the relevant comparison is against the complete complement. In this case, the evidence slightly favors the complete joint complement over the joint predicted hypothesis, $\mathrm{BF}_{1,c^{*}}^{\text{(joint)}}=0.36$. That is, the data are somewhat more likely if the predictions fail in at least one study than if they hold throughout. Notably, testing against the joint unconstrained hypothesis yields a virtually identical value, $\mathrm{BF}_{1,u}^{\text{(joint)}}=0.36$.

\subsection{Network control effects on trust and cooperation %\texorpdfstring{\parencite{volkerFutureMadeToday2022}}{(Volker, 2022)}
}

\textcite{volkerFutureMadeToday2022} investigated the effects of dyadic and network embeddedness on trust and cooperation. \textit{Dyadic embeddedness} occurs when a person X can draw upon previous experiences with person Y to guide the decision of how much---or whether---to trust them in a current exchange. In turn, \textit{network embeddedness} occurs when person X can access information from third parties about their previous interactions with person Y. The author tested the overarching hypothesis that embeddedness promotes trust and cooperation by synthesizing evidence across eight heterogeneous studies. For example, these studies differed in the type of embeddedness that was assessed (dyadic vs.\ network), the trust game that was used (e.g., prisoner's dilemma, investment game), and the outcome measure that was considered (trustfulness, trustworthiness, or cooperation).

For each study, \textcite{volkerFutureMadeToday2022} computed all $\mathrm{BF}_{1,u}^{(t)}$ and $\mathrm{BF}_{1,c}^{(t)}$ and combined them using BES to obtain $\mathrm{BF}_{1,u}^{(\text{joint})}$ and $\mathrm{BF}_{1,c}^{(\text{joint})}$. Table~\ref{tab:volker} presents the study-specific Bayes factors reported by \textcite{volkerFutureMadeToday2022}, together with the joint Bayes factors that we computed from them.%
    \footnote{Joint Bayes factors were computed from the three-decimal values in Table~4 of \textcite{volkerFutureMadeToday2022}, the most precise values available. Because of this rounding, our $\mathrm{BF}_{1,u}^{(\text{joint})}$ differs by 0.03 from the originally reported value. We computed $\mathrm{BF}_{1,c^{*}}^{(\text{joint})}$ via Eq.~\eqref{eq:bf-cstar-u-joint}, which uses the $\mathrm{BF}_{1,u}^{(t)}$. Computing it from the separately rounded $\mathrm{BF}_{1,c}^{(t)}$ via Eq.~\eqref{eq:bf-1-cstar-equal} instead gives 7.28, but for exact inputs, both computations yield equivalent results.}
As in the previous example, testing against the incomplete joint complement (i.e., $H_1^{(t)}$ is false in \textit{all} studies) provides overwhelming evidence in favor of the joint predicted hypothesis, $\mathrm{BF}_{1,c}^{(\text{joint})} = 5.23 \times 10^{11}$. Testing against the complete joint complement likewise supports the joint predicted hypothesis, but with substantially weaker evidence, $\mathrm{BF}_{1,c^{*}}^{(\text{joint})} = 7.24$. Again, we observe that the Bayes factor testing against the complete complement is virtually identical to that obtained from testing against the joint unconstrained hypothesis, $\mathrm{BF}_{1,u}^{(\text{joint})} = 7.24$.

    \begin{table}[!tb]
    \captionsetup{justification=raggedright, singlelinecheck=false}
    \caption{Study-specific and joint Bayes factors testing whether network embeddedness promotes trust and cooperation \parencite{volkerFutureMadeToday2022}}
    \label{tab:volker}

    \centering
    \begin{tabular}{l>{\centering\arraybackslash}c>{\centering\arraybackslash}c>{\centering\arraybackslash}c}
    \hline
    Study &
    Outcome &
    $\mathrm{BF}_{1,u}$ &
    $\mathrm{BF}_{1,c}$ \\
    \hline

    \multirow[t]{2}{*}{Bolton et al. (2004)}
        & Trustfulness    & 1.47 & 2.80 \\
        & Trustworthiness & 0.64 & 0.47 \\

    \multirow[t]{2}{*}{Duffy et al. (2013): No-NetMin}
        & Trustfulness    & 1.29 & 1.83 \\
        & Trustworthiness & 1.82 & 10.32 \\

    \multirow[t]{2}{*}{Duffy et al. (2013): No-NetFull}
        & Trustfulness    & 1.80 & 8.96  \\
        & Trustworthiness & 2.00 & $2.60 \times 10^{4}$  \\

    \multirow[t]{2}{*}{Buskens et al. (2010)}
        & Trustfulness    & 1.85 & 12.17 \\
        & Trustworthiness & 1.68 & 5.25  \\

    \multirow[t]{2}{*}{Van Miltenburg et al. (2012)}
        & Trustfulness    & 0.67 & 0.51  \\
        & Trustworthiness & 1.11 & 1.26  \\

    \multirow[t]{2}{*}{Frey et al. (2019)}
        & Trustfulness    & 1.43 & 2.49  \\
        & Trustworthiness & 2.00 & $1.01 \times 10^{3}$ \\

    \multirow[t]{2}{*}{Barrera \& Buskens (2009)}
        & Trustfulness & 0.48 & 0.32 \\
        & Trustworthiness & 1.45  & 2.63 \\

    Seinen \& Schram (2006)
        & Cooperation & 1.90 & 19.35 \\

    Corten et al. (2016)
        & Cooperation & 0.10 & 0.05 \\

   \hline
    BES: product BF
        &
        & 7.24
        & $5.23 \times 10^{11 \, \dagger}$ \\
    BES: complete complement
        &
        &
        & $7.24^{*}$ \\
    \hline
    \end{tabular}
    \vspace{2mm}
    \begin{flushleft}
    \footnotesize

   \textit{Note.} Study-specific hypotheses: $H_1: \beta>0$, $H_u: \beta\in\mathbb{R}$, and $H_c: \beta\le 0$. Following \textcite[p.~32]{volkerFutureMadeToday2022}, $\beta$ has a normal prior $\mathcal{N}(0, \hat{\Sigma}_{\beta}/b)$ under $H_u$; hence, $c_1 = 0.50$ in all studies. $^{\dagger}$Evidence against the incomplete joint complement. $^{*}$Evidence against the complete joint complement. Study-specific Bayes factors are taken from Table~2 in \textcite{volkerFutureMadeToday2022}, which reports them rounded to two decimals. Trustfulness and trustworthiness BFs from the same study are based on the same participants, so the independence assumption of BES is questionable for these pairs; we nevertheless follow the specification of \textcite{volkerFutureMadeToday2022}.
    \end{flushleft}
    \end{table}


\section{Discussion}

We showed that the conclusions drawn from Bayesian evidence synthesis \parencite{kuiperCombiningStatisticalEvidence2013, klugkistBayesianEvidenceSynthesis2025} depend on the joint alternative hypothesis across studies against which the theory is evaluated. We considered four possible joint alternative hypotheses: the null, the unconstrained, the incomplete complement, and the complete complement. Among these, we argued that the incomplete joint complement is of limited relevance for theory testing because it rarely represents a substantively interesting claim.

\subsection{The incomplete vs. the complete joint complement}
\label{sec:complete-incomplete}

Multiplying study-specific Bayes factors that test theoretically derived order constraints against their complements implicitly uses the incomplete joint complement as the alternative hypothesis, which states that the theoretical prediction is wrong in every study. For theory testing, a more appropriate alternative is the \textit{complete} joint complement, which states that the prediction is wrong in at least one study, that is, the theory does not hold across all studies. We showed how to compute the Bayes factor against the complete joint complement from the study-specific Bayes factors alone. In our empirical examples, the two alternatives led to substantially different, in one case even opposite, conclusions about the theory.

The discrepancy between the two types of complements is reflected in the geometry of the joint parameter space. For more than one study, the incomplete complement is a strict subset of the complete complement. As studies are added, the former shrinks exponentially, whereas the latter approaches the entire parameter space. In Example A (Section~\ref{sec:example-a-choice-probabilities}), the two complements coincide for a single study, each covering 50\% of the parameter space; for two studies, they cover 25\% and 75\%, and for three studies, 12.5\% and 87.5\%, respectively. Put differently, the complete complement has a higher complexity, meaning that the parameters are a priori more likely to be in line with the encompassing prior \parencite{klugkistBayesFactorInequality2007}.

\subsection{Choosing a competing joint hypothesis}

Setting aside the incomplete joint complement, we argue that there is no universally correct choice among the three remaining joint alternatives. Which alternative hypothesis is appropriate depends on the scientific question.

The joint null is appropriate when the question is whether an effect is present in all studies or absent in all studies, for example, in clinical trials, where it provides a clear benchmark for concluding that a treatment has no effect. However, neither the joint null nor the joint predicted hypothesis covers more complex scenarios in which the effect is present in some studies but absent in others. In such cases, $\mathrm{BF}_{1,0}^{(\text{joint})}$ necessarily contrasts two false hypotheses, so whichever hypothesis the Bayes factor favors, the conclusion is misleading. Moreover, because Bayes factors based on commonly used alternative priors accumulate evidence against a false point null much faster than evidence for a true one \parencite{johnsonUseNonLocalPrior2010}, the few studies in which the effect is present can provide overwhelming support for the predictions that outweighs the comparatively weak evidence for the null from the remaining studies.

Unlike the joint null, the complete complement explicitly includes such mixed patterns, in which the predictions hold in some studies but fail in others. It is therefore the natural choice for exposing a theory to falsification attempts. If each study is an independent test of the theory, the theory implies that its predictions hold in all studies, and so a violation in any single study counts against it. The complete complement thus asks whether the theory makes \textit{any} incorrect prediction. In contrast, the joint unconstrained hypothesis does not pit the theory against its logical complement, but asks whether imposing the theoretical constraints is justified at all, which is the rationale underlying the encompassing Bayes factor \parencite{klugkistBayesFactorInequality2007, hoijtinkInformativeHypotheses2012}.

When testing order constraints, the unconstrained hypothesis has often been used as an alternative hypothesis \parencite[e.g.,][]{heckMultinomialModelsLinear2019, zwilling2019qtest, hoijtinkInformativeHypotheses2012}. However, using the unconstrained, encompassing model has been criticized because its predictions overlap with the constrained predicted hypothesis \parencite{shiffrinBayesFactorsRelations2016, chandramouliCommentaryGronauWagenmakers2019}. Two objections follow on this basis, but we argue that neither disqualifies $H_u^{(\text{joint})}$ in the context of theory testing.

First, the comparison of the competing models is said to lack a clear, qualitative contrast: Whether a parameter falls above or below a certain boundary cannot be answered for values in the shared region and is better addressed by comparing disjoint regions of the parameter space \parencite{shiffrinBayesFactorsRelations2016}. We agree that directional questions require disjoint hypotheses, which in BES is achieved by using the complete complement. However, the comparison against $H_u^{(\text{joint})}$ asks a different question, namely, whether committing to the theoretical constraints is warranted relative to remaining agnostic, and these positions do differ qualitatively. The overlap between the predicted and the unconstrained hypothesis is not incidental but constitutive: The comparison is informative precisely because $H_1$ is a restriction of $H_u$. Using encompassing, saturated models as ``catch-all'' alternatives has also been shown to be beneficial in cognitive modeling \parencite[e.g.,][]{heck2017information, lee2016bayesian}.

Second, because $H_1$ is nested in $H_u$, the attainable evidence in favor of the predictions is necessarily bounded. Within a single study, $\mathrm{BF}_{1,u}^{(t)}$ cannot exceed $1/c_1^{(t)}$, the inverse complexity of the predicted hypothesis, no matter how much data are collected or how strongly the results agree with the predicted order constraints \parencite{heckTestingOrderConstraints2015, klugkistBayesianEvidenceSynthesis2025}. In contrast, $\mathrm{BF}_{1,c}^{(t)}$ is unbounded, so the two Bayes factors can diverge substantially when the data are in line with the order constraints \parencite{shiffrinBayesFactorsRelations2016}. For a single study, this bound can indeed be restrictive. In BES, however, the relevant bound is $1/c_1^{(\text{joint})}$, and because the joint complexity $c_1^{(\text{joint})} = \prod_t c_1^{(t)}$ shrinks exponentially with the number of studies (cf. the geometric illustration of the joint parameter space in Section~\ref{sec:complete-incomplete}), the bound grows accordingly. If $c_1^{(t)} \le 0.50$ in each study, which will be the case if the predictions are sufficiently precise, the bound is at least $2^T$, so it exceeds 10 with four studies and 100 with seven. The bound thus quickly ceases to be a practical constraint as studies accumulate.%
    \footnote{For less precise predicted hypotheses ($c_1^{(t)}>0.50$), the bound grows more slowly. For example, with $c_1^{(t)}=0.90$ in all studies, it exceeds 10 only with 22 studies and 100 only with 44.}

Does the choice between $H_u^{(\text{joint})}$ and $H_{c^*}^{(\text{joint})}$ matter in practice? The two Bayes factors were virtually identical in our empirical examples, and because the alternatives cover increasingly similar regions of the parameter space as studies are added, one may expect the Bayes factors to converge. However, their ratio, $\mathrm{BF}_{1,u}^{(\text{joint})}/\mathrm{BF}_{1,c^*}^{(\text{joint})} = \mathrm{BF}_{c^*,u}^{(\text{joint})}$, need not approach one. Whether the two Bayes factors converge depends on the weighting of the components of the complete complement (i.e., the a priori plausibility of different patterns of correct and incorrect predictions across studies) and on the evidence provided by the studies (for detailed derivations, see Section~S1 of the Supplementary Material, \suppurl). Under complexity weighting, the two Bayes factors always point in the same direction, with $\mathrm{BF}_{1,c^*}^{(\text{joint})}$ at least as extreme as $\mathrm{BF}_{1,u}^{(\text{joint})}$. Moreover, their ratio approaches one as studies are added, unless the posterior concentrates in the joint predicted hypothesis. Under equal weighting, which we consider more appropriate for theory testing (Section~\ref{sec:step-3-compute-bf-cstar-u-joint}), neither property holds in general. Depending on the complexities and the evidence, the two Bayes factors can differ by several orders of magnitude even when many studies are synthesized, and may even point in opposite directions. This explains their agreement in our empirical examples. All predicted hypotheses had the same complexity of 50\% (i.e., $c_1^{(t)} = 0.50$), so the two prior weightings coincide, and the posterior probability that all predictions hold was very small (about $0.30\%$ and $0.01\%$). Hence, the choice between $H_u^{(\text{joint})}$ and $H_{c^*}^{(\text{joint})}$ should rest on the substantive question rather than on the expectation that they will agree. If it is unclear which alternative is more appropriate, both can be reported.

\subsection{Heterogeneity across studies}

Irrespective of which alternative hypothesis is used, the joint Bayes factor can obscure heterogeneity in the evidence across studies. In the towel-reuse example \parencite{scheibehenneBayesianEvidenceSynthesis2016}, \mbox{$\mathrm{BF}_{1,u}^{(\text{joint})} = 0.36$} indicates weak evidence against the theory, yet the study-specific $\mathrm{BF}_{1,u}^{(t)}$ range from 0.04 (strong evidence against the prediction) to 1.97 (weak evidence in favor, close to the attainable maximum of 2). The contrast is even starker for the complement. Although the study-specific $\mathrm{BF}_{1,c}^{(t)}$ range from 0.02 to 75.34, the joint Bayes factor is \mbox{$\mathrm{BF}_{1,c^{*}}^{(\text{joint})} = 0.36$}. Likewise, in the network-embeddedness example, the $\mathrm{BF}_{1,c}^{(t)}$ range from 0.05 to $2.60 \times 10^4$, whereas \mbox{$\mathrm{BF}_{1,c^{*}}^{(\text{joint})} = 7.24$} indicates moderate evidence. The joint Bayes factor can thus mask effect heterogeneity that shows up as strong evidence in both directions at the study level. Because the product depends on the magnitudes of the study-specific Bayes factors, not only their direction, the joint Bayes factor need not favor the hypothesis supported by most studies. This is exactly what happened in the towel-reuse example: Both joint Bayes factors favored their respective alternatives, although only two of the seven experiments did.

These problems reflect a general limitation of any single summary measure of evidence, not an argument for one joint alternative over another. As a remedy, researchers should inspect the distribution of study-specific Bayes factors alongside the joint Bayes factor, for example, in a histogram of the study-specific $\log \mathrm{BF}$ values \parencite[for examples with heterogeneous person-specific Bayes factors, see][]{stephanBayesianModelSelection2009}. Beyond inspecting the distribution, researchers can plot the study-specific Bayes factors against study characteristics to reveal systematic patterns that may point to potential moderators.

\subsection{Conclusion}

We have shown that researchers should be careful when multiplying Bayes factors across studies. By doing so, they may inadvertently test their theory against a competing hypothesis that does not reflect a theoretically relevant position. As always when using Bayes factors, the specification of the hypotheses being tested should be considered carefully \parencite{heckReviewApplicationsBayes2023, tendeiroReviewIssuesNull2019}. In Bayesian evidence synthesis, it is sensible to test one's predictions not only against the joint null hypothesis, but also against hypotheses that allow for heterogeneity across studies, such as the complete complement or the unconstrained hypothesis. Whichever alternative is chosen, the joint Bayes factor should be reported alongside the study-specific Bayes factors from which it was computed, so that readers can judge whether the aggregate reflects a consistent pattern of evidence.
%-------------------------------------------------------------------------------------------------------
\FloatBarrier % force LaTeX to place all pending figures before

%--------------------------
%           Credits       %
%--------------------------
\section*{Acknowledgments \& Disclosures}

Preprint version 1.0. We used Claude \parencite{anthropicClaudeOpus2026} for language improvement, for producing Python code for the plots, for reviewing R code of the computations, and for derivations in the Appendix and Supplementary Material. We declare that there are no conflicts of interest with respect to the authorship or the publication of this article.

\section*{Author Contributions}

Conceptualization: L.G., D.W.H.
Data Curation: L.G.
Formal Analysis: L.G., D.W.H.
Supervision: D.W.H.
Visualization: L.G. 
Writing--Original Draft: L.G.
Writing--Review \& Editing: L.G., D.W.H.

\printbibliography

\appendix
\section{Deriving expressions for \texorpdfstring{\boldmath$\mathrm{BF}^{(\mathrm{joint})}_{1,c^*}$}{BF1c*}}
\label{app:closed-form}

In the following, we derive simplified expressions for $\mathrm{BF}^{(\text{joint})}_{1,c^*}$
under both prior weightings for the complete complement, which can be computed without enumerating the $2^T-1$ components. Each derivation first simplifies Eq.~\eqref{eq:bf-cstar-u-joint}
for $\mathrm{BF}^{(\text{joint})}_{c^*,u}$ and then obtains
$\mathrm{BF}^{(\text{joint})}_{1,c^*}$ via the transitivity of Bayes factors. The Bayes factor comparing two hypotheses equals the ratio of their Bayes factors against a common reference hypothesis, here the joint unconstrained hypothesis, whose marginal likelihood cancels:
\begin{flalign}
\derivindent\mathrm{BF}^{(\text{joint})}_{1,c^*}
&= \frac{P\bigl(D \mid H_1^{(\text{joint})}\bigr)}
        {P\bigl(D \mid H_{c^*}^{(\text{joint})}\bigr)} && \nonumber\\
&= \frac{P\bigl(D \mid H_1^{(\text{joint})}\bigr)\big/P\bigl(D \mid H_u^{(\text{joint})}\bigr)}
        {P\bigl(D \mid H_{c^*}^{(\text{joint})}\bigr)\big/P\bigl(D \mid H_u^{(\text{joint})}\bigr)} \nonumber\\
&= \frac{\mathrm{BF}^{(\text{joint})}_{1,u}}{\mathrm{BF}^{(\text{joint})}_{c^*,u}}.
   \label{eq:transitivity}
\end{flalign}
Both derivations also use the encompassing prior approach in Eqs.~\eqref{eq:bf-1u-t} and~\eqref{eq:bf-cu-t}, which
express each study-specific Bayes factor as the ratio of model fit and complexity,
$\mathrm{BF}^{(t)}_{1,u} = f_1^{(t)}/c_1^{(t)}$ and
$\mathrm{BF}^{(t)}_{c,u} = \bigl(1-f_1^{(t)}\bigr)/\bigl(1-c_1^{(t)}\bigr)$.

\subsection{Complexity weighting}
\label{app:complexity}

Under the independence assumptions discussed in Section~\ref{sec:joint-predicted}, the
joint posterior factorizes across studies. The fit of component $H_\Omega$, that is, its
posterior probability, is therefore
$f_\Omega = \prod_{t\in\Omega}\bigl(1-f_1^{(t)}\bigr)\prod_{t\notin\Omega} f_1^{(t)}$.
We further define the joint complexity as
$c_1^{(\text{joint})} = \prod_{t=1}^{T} c_1^{(t)}$ and, analogously, the joint fit as
$f_1^{(\text{joint})} = \prod_{t=1}^{T} f_1^{(t)}$, the fit of the joint predicted hypothesis ($\Omega = \emptyset$). Because the $2^T$ components partition the joint parameter space, their fits sum to one, so the fit of the complete complement, which comprises all components except $\Omega = \emptyset$, is
\begin{equation}
\label{eq:fit-cstar}
\sum_{\emptyset\neq\Omega} f_\Omega = 1 - f_1^{(\text{joint})}
= 1 - \prod_{t=1}^{T} f_1^{(t)}.
\end{equation}
When each component's Bayes factor is expressed as fit over complexity, the complexity weights cancel the component complexities:
\begin{flalign}
\derivindent\mathrm{BF}^{(\text{joint})}_{c^*,u}
&= \sum_{\emptyset\neq\Omega\subseteq\{1,\dots,T\}} \pi_\Omega
\prod_{t\in\Omega}\mathrm{BF}^{(t)}_{c,u}\prod_{t\notin\Omega}\mathrm{BF}^{(t)}_{1,u}
   && \text{Eq.~\eqref{eq:bf-cstar-u-joint}} \nonumber\\
&= \sum_{\emptyset\neq\Omega} \pi_\Omega
   \prod_{t\in\Omega}\frac{1-f_1^{(t)}}{1-c_1^{(t)}}
   \prod_{t\notin\Omega}\frac{f_1^{(t)}}{c_1^{(t)}}
   && \text{Eqs.~\eqref{eq:bf-1u-t} and \eqref{eq:bf-cu-t}} \nonumber\\
&= \sum_{\emptyset\neq\Omega} \pi_\Omega\,\frac{f_\Omega}{c_\Omega}
   && \text{combine numerators and denominators} \nonumber\\
&= \sum_{\emptyset\neq\Omega} \frac{c_\Omega}{1-\prod_{t} c_1^{(t)}}
   \cdot\frac{f_\Omega}{c_\Omega}
   && \text{substitute complexity weights} \nonumber\\
&= \frac{1}{1-\prod_{t} c_1^{(t)}}\sum_{\emptyset\neq\Omega} f_\Omega
   && \text{$c_\Omega$ cancels; factor out the constant} \nonumber\\
&= \frac{1-\prod_{t} f_1^{(t)}}{1-\prod_{t} c_1^{(t)}}
   && \text{Eq.~\eqref{eq:fit-cstar}} \nonumber\\
&= \frac{1-c_1^{(\text{joint})}\,\mathrm{BF}^{(\text{joint})}_{1,u}}
        {1-c_1^{(\text{joint})}}
   && \text{Eq.~\eqref{eq:bf-1u-t}: } f_1^{(t)} = c_1^{(t)}\,\mathrm{BF}^{(t)}_{1,u}.
   \label{eq:bf-cstar-u-complexity}
\end{flalign}
The final expression has the same form as Eq.~\eqref{eq:bf-cu-t}, with the study-specific quantities replaced by their joint counterparts. Substituting Eq.~\eqref{eq:bf-cstar-u-complexity} into Eq.~\eqref{eq:transitivity}
yields Eq.~\eqref{eq:bf-1-cstar-complexity}:
\begin{flalign*}
\derivindent\mathrm{BF}^{(\text{joint})}_{1,c^*}
&= \frac{\mathrm{BF}^{(\text{joint})}_{1,u}}
        {\dfrac{1-c_1^{(\text{joint})}\,\mathrm{BF}^{(\text{joint})}_{1,u}}
               {1-c_1^{(\text{joint})}}}
   && \text{Eq.~\eqref{eq:bf-cstar-u-complexity}} \\
&= \frac{\bigl(1-c_1^{(\text{joint})}\bigr)\,\mathrm{BF}^{(\text{joint})}_{1,u}}
        {1-c_1^{(\text{joint})}\,\mathrm{BF}^{(\text{joint})}_{1,u}}
   && \text{multiply numerator and denominator by } 1-c_1^{(\text{joint})} \\
&= \frac{\bigl(1-\prod_{t=1}^{T} c_1^{(t)}\bigr)\prod_{t=1}^{T}\mathrm{BF}^{(t)}_{1,u}}
        {1-\prod_{t=1}^{T}\bigl(c_1^{(t)}\,\mathrm{BF}^{(t)}_{1,u}\bigr)}
   && \text{write joint quantities as products}.
\end{flalign*}
% Substituting $\mathrm{BF}^{(\text{joint})}_{1,u} = f_1^{(\text{joint})}/c_1^{(\text{joint})}$
% into the second line shows that this Bayes factor is the posterior odds divided by the
% prior odds of the joint predicted hypothesis against its complement,
% \begin{equation*}
% \mathrm{BF}^{(\text{joint})}_{1,c^*}
% = \frac{f_1^{(\text{joint})}\big/\bigl(1-f_1^{(\text{joint})}\bigr)}
%        {c_1^{(\text{joint})}\big/\bigl(1-c_1^{(\text{joint})}\bigr)}.
% \end{equation*}
% Unlike Eq.~\eqref{eq:bf-1-cstar-equal}, this expression is not a function of the
% study-specific Bayes factors $\mathrm{BF}^{(t)}_{1,c}$ alone.

\subsection{Equal weighting}
\label{app:equal}

Under equal weighting, the weights do not depend on the complexities $c_\Omega$. Hence,
unlike in Appendix~\ref{app:complexity}, they do not cancel the complexities in the
denominators of the components' Bayes factors, and the derivation proceeds differently.
It uses the fact that expanding the product
$\prod_{t}\bigl(\mathrm{BF}^{(t)}_{1,u}+\mathrm{BF}^{(t)}_{c,u}\bigr)$
yields one term for each of the $2^T$ components, namely the product of
$\mathrm{BF}^{(t)}_{c,u}$ for the studies in $\Omega$ and $\mathrm{BF}^{(t)}_{1,u}$ for
the remaining studies:
\begin{equation}
\label{eq:expand-product}
\sum_{\Omega\subseteq\{1,\dots,T\}}
\prod_{t\in\Omega}\mathrm{BF}^{(t)}_{c,u}\prod_{t\notin\Omega}\mathrm{BF}^{(t)}_{1,u}
= \prod_{t=1}^{T}\bigl(\mathrm{BF}^{(t)}_{1,u}+\mathrm{BF}^{(t)}_{c,u}\bigr).
\end{equation}
With equal weights, the factor $1/(2^T-1)$ can be pulled out of the sum, and the sum over the non-empty components equals the full sum in Eq.~\eqref{eq:expand-product} minus the term for $\Omega = \emptyset$, which is $\mathrm{BF}^{(\text{joint})}_{1,u}$:
\begin{flalign}
\derivindent\mathrm{BF}^{(\text{joint})}_{c^*,u}
&= \sum_{\emptyset\neq\Omega\subseteq\{1,\dots,T\}} \pi_\Omega
   \prod_{t\in\Omega}\mathrm{BF}^{(t)}_{c,u}\prod_{t\notin\Omega}\mathrm{BF}^{(t)}_{1,u}
   && \text{Eq.~\eqref{eq:bf-cstar-u-joint}} \nonumber\\
&= \frac{1}{2^T-1}\sum_{\emptyset\neq\Omega}
   \prod_{t\in\Omega}\mathrm{BF}^{(t)}_{c,u}\prod_{t\notin\Omega}\mathrm{BF}^{(t)}_{1,u}
   && \text{substitute equal weights} \nonumber\\
&= \frac{1}{2^T-1}\Biggl[\sum_{\Omega\subseteq\{1,\dots,T\}}
   \prod_{t\in\Omega}\mathrm{BF}^{(t)}_{c,u}\prod_{t\notin\Omega}\mathrm{BF}^{(t)}_{1,u}
   - \prod_{t=1}^{T}\mathrm{BF}^{(t)}_{1,u}\Biggr]
   && \text{add and subtract } \Omega=\emptyset \nonumber\\
&= \frac{\prod_{t=1}^{T}\bigl(\mathrm{BF}^{(t)}_{1,u}+\mathrm{BF}^{(t)}_{c,u}\bigr)
   -\mathrm{BF}^{(\text{joint})}_{1,u}}{2^T-1}
   && \text{Eq.~\eqref{eq:expand-product}}. \label{eq:bf-cstar-u-equal}
\end{flalign}
Substituting Eq.~\eqref{eq:bf-cstar-u-equal} into Eq.~\eqref{eq:transitivity} yields
Eq.~\eqref{eq:bf-1-cstar-equal}:
\begin{flalign*}
\derivindent\mathrm{BF}^{(\text{joint})}_{1,c^*}
&= \frac{\mathrm{BF}^{(\text{joint})}_{1,u}}
        {\dfrac{\prod_{t}\bigl(\mathrm{BF}^{(t)}_{1,u}+\mathrm{BF}^{(t)}_{c,u}\bigr)
        -\mathrm{BF}^{(\text{joint})}_{1,u}}{2^T-1}}
   && \text{Eq.~\eqref{eq:bf-cstar-u-equal}} \\
&= \frac{\bigl(2^T-1\bigr)\prod_{t}\mathrm{BF}^{(t)}_{1,u}}
        {\prod_{t}\bigl(\mathrm{BF}^{(t)}_{1,u}+\mathrm{BF}^{(t)}_{c,u}\bigr)
        -\prod_{t}\mathrm{BF}^{(t)}_{1,u}}
   && \text{multiply numerator and denominator by } 2^T-1 \\
&= \frac{2^T-1}
        {\prod_{t}\bigl(1+\mathrm{BF}^{(t)}_{c,u}/\mathrm{BF}^{(t)}_{1,u}\bigr)-1}
   && \text{divide numerator and denominator by }
      \textstyle\prod_{t}\mathrm{BF}^{(t)}_{1,u} \\
&= \frac{2^T-1}{\prod_{t=1}^{T}\bigl(1+1/\mathrm{BF}^{(t)}_{1,c}\bigr)-1}
   && \text{substitute } \mathrm{BF}^{(t)}_{c,u}/\mathrm{BF}^{(t)}_{1,u} = 1/\mathrm{BF}^{(t)}_{1,c}.
\end{flalign*}
Unlike Eq.~\eqref{eq:bf-1-cstar-complexity}, which requires the study-specific Bayes
factors $\mathrm{BF}^{(t)}_{1,u}$ together with the complexities $c_1^{(t)}$, this
expression requires only the study-specific Bayes factors $\mathrm{BF}^{(t)}_{1,c}$.

% \subsection{Adding an uninformative study}
% \label{app:uninformative}

% Suppose that a study $T+1$ with
% $\mathrm{BF}^{(T+1)}_{1,u} = \mathrm{BF}^{(T+1)}_{c,u} = 1$, that is,
% $f_1^{(T+1)} = c_1^{(T+1)}$, is added. Each previous non-empty component $\Omega$ then
% appears twice, once with $H_1$ and once with $H_c$ in the new study. Both versions have
% the same Bayes factor as component $\Omega$ had before, and their combined weight is
% proportional to the previous weight of component $\Omega$: under equal weighting, each
% version receives weight $1/(2^{T+1}-1)$; under complexity weighting, their prior masses
% $c_\Omega\, c_1^{(T+1)}$ and $c_\Omega\bigl(1-c_1^{(T+1)}\bigr)$ sum to $c_\Omega$.
% Together, these components therefore reproduce the previous
% $\mathrm{BF}^{(\text{joint})}_{c^*,u}$. The only new component is the one in which the
% theory fails only in the new study; since the predicted hypothesis holds in all previous
% studies, its Bayes factor is $\mathrm{BF}^{(\text{joint})}_{1,u}$. The updated Bayes
% factor for the complete complement is therefore
% \begin{equation}
% \label{eq:uninformative-mixture}
% \mathrm{BF}^{(\text{joint}),\,\text{new}}_{c^*,u}
% = w\,\mathrm{BF}^{(\text{joint})}_{1,u} + (1-w)\,\mathrm{BF}^{(\text{joint})}_{c^*,u},
% \end{equation}
% where $w$ is the weight of the new component,
% \begin{equation*}
% w = \frac{1}{2^{T+1}-1} \quad \text{(equal weighting)},
% \qquad
% w = \frac{\bigl(1-c_1^{(T+1)}\bigr)\,c_1^{(\text{joint})}}
%          {1-c_1^{(T+1)}\,c_1^{(\text{joint})}} \quad \text{(complexity weighting)}.
% \end{equation*}
% Because $\mathrm{BF}^{(\text{joint})}_{1,u}$ is unchanged, the updated
% $\mathrm{BF}^{(\text{joint})}_{1,c^*}
% = \mathrm{BF}^{(\text{joint})}_{1,u}\big/\mathrm{BF}^{(\text{joint}),\,\text{new}}_{c^*,u}$
% lies between its previous value and one. Under equal weighting, $w$ depends only on the
% number of studies. Under complexity weighting, $w$ increases as $c_1^{(T+1)}$ decreases.
% Under both weightings, $w$ decreases rapidly as studies accumulate.
% Table~\ref{tab:uninformative} illustrates the computation for Example A.

% \begin{table}[ht]
% \centering
% \caption{Effect of adding an uninformative fourth study to Example A}
% \label{tab:uninformative}
% \begin{tabular}{lcccc}
% \toprule
% Weighting & $c_1^{(4)}$ & $w$
%   & $\mathrm{BF}^{(\text{joint}),\,\text{new}}_{c^*,u}$
%   & $\mathrm{BF}^{(\text{joint}),\,\text{new}}_{1,c^*}$ \\
% \midrule
% Equal      & any & 0.067 & 0.941 & 2.01 \\
% Complexity & 0.1 & 0.114 & 0.989 & 1.91 \\
% Complexity & 0.5 & 0.067 & 0.941 & 2.01 \\
% Complexity & 0.9 & 0.014 & 0.887 & 2.13 \\
% \bottomrule
% \end{tabular}

% \medskip
% \begin{minipage}{\linewidth}
% \small
% \textit{Note.} Before adding the study,
% $\mathrm{BF}^{(\text{joint})}_{1,u} = 1.89$,
% $\mathrm{BF}^{(\text{joint})}_{c^*,u} = 0.873$, and
% $\mathrm{BF}^{(\text{joint})}_{1,c^*} = 2.17$. The new study has
% $\mathrm{BF}^{(4)}_{1,u} = \mathrm{BF}^{(4)}_{c,u} = 1$, so that
% $\mathrm{BF}^{(\text{joint})}_{1,u}$ remains 1.89. The weight $w$ is
% $1/(2^4-1)$ under equal weighting and follows from $c_1^{(\text{joint})} = 0.5^3 = 0.125$
% under complexity weighting. The updated values follow from
% Eq.~\eqref{eq:uninformative-mixture} and
% $\mathrm{BF}^{(\text{joint}),\,\text{new}}_{1,c^*}
% = 1.89/\mathrm{BF}^{(\text{joint}),\,\text{new}}_{c^*,u}$.
% \end{minipage}
% \end{table}

\end{document}